\begin{table}[tp]
\centering
\caption{Inexact gradients for strongly convex optimization: \href{https://github.com/k3nfalt/longtails/blob/main/LongTails/Setups/InexactGradientStronglyConvex.lean}{Inexact\allowbreak{}Gradient\allowbreak{}Strongly\allowbreak{}Convex}.}\label{tab:InexactGradientStronglyConvex}
\parbox{\linewidth}{\footnotesize On a real Hilbert space, differentiable $f$ satisfies the quadratic upper and lower gradient models with constants $L$ and $\mu$, where $0<\mu\le L$. A minimizer $x_\star$ satisfies $\|x_0-x_\star\|\le R$, $R\ge0$. Every oracle reply obeys $\|g-\nabla f(x)\|\le\alpha\|\nabla f(x)\|+\delta$, with $0\le\alpha<1$ and $\delta\ge0$; errors may depend on the call and seed. Target: $f(x_N)-f(x_\star)\le\varepsilon$ for every seed, with $\varepsilon>0$ and $136\delta^2\le\mu(1-\alpha)^2\varepsilon$.}\par\smallskip
\begin{threeparttable}
\footnotesize
\setlength{\tabcolsep}{4pt}\renewcommand{\arraystretch}{1.25}
\begin{tabularx}{\linewidth}{@{}>{\raggedright\arraybackslash}p{0.20\linewidth}>{\raggedright\arraybackslash}p{0.30\linewidth}>{\raggedright\arraybackslash}p{0.28\linewidth}>{\raggedright\arraybackslash}X@{}}
\toprule
\textbf{Result} & \textbf{Oracle complexity} & \textbf{Novelty} & \textbf{Limitations} \\
\midrule
\textbf{(1)} \href{https://github.com/k3nfalt/longtails/blob/main/LongTails/Theorems/AcceleratedInexactStronglyConvex.lean}{Accelerated\allowbreak{}Inexact\allowbreak{}Strongly\allowbreak{}Convex} \textbf{(Ours)} $(\star\star\star)$ \href{https://github.com/k3nfalt/longtails/blob/main/scripts/reviews/AcceleratedInexactStronglyConvex.md}{(review)}
& $\mathcal{O}\!\left(1+\left[\sqrt{\frac{\kappa}{1-\alpha}}+\frac{\alpha^2\kappa}{1-\alpha}\right]\ell\right)$
& --- & --- \\
\midrule[1.2pt]
RE-AGM \citep[Thm.~6.1]{vasinGasnikov2026CompositeV1}\par $\alpha\le1/3$, $E=\mathbb{R}^d$
& $\mathcal{O}\!\left(1+(\sqrt\kappa+\alpha\kappa)\ell\right)$\par $\varepsilon\gtrsim\min\{\sqrt\kappa,\alpha^{-1}\}\delta^2/\mu$
& i) \textbf{(1)}: gain $\min\{\alpha\sqrt\kappa,\alpha^{-1}\}$ for $\kappa^{-1/2}\le\alpha\le1/3$.\newline ii) \textbf{(1)}: certified threshold $\delta^2/\mu$ instead of $\min\{\sqrt\kappa,\alpha^{-1}\}\delta^2/\mu$ on this range.\newline iii) \textbf{(1)}: matches rate for $\alpha\lesssim\kappa^{-1/2}$.
& --- \\
\midrule
RE-AGM \citep[Thms.~4--5]{vasinEtAl2025RelativeV3}\par Special case: $\delta=0$, $E=\mathbb{R}^d$
& $\mathcal{O}\!\left(1+(\sqrt\kappa+\alpha\kappa)\ell\right)$\par $0\le\alpha\le1/3$
& i) \textbf{(1)}: $\sqrt\kappa$ dependence retained up to $\alpha\asymp\kappa^{-1/4}$ instead of $\kappa^{-1/2}$.\newline ii) \textbf{(1)}: $\alpha^2\kappa$ instead of $\alpha\kappa$ for $\kappa^{-1/4}\le\alpha\le1/3$.\newline iii) \textbf{(1)}: scope: permits $\delta>0$; matched rate for $\alpha\lesssim\kappa^{-1/2}$.
& --- \\
\midrule
Fast gradient \citep[Thm.~3.3]{gaucherEtAl2026InexactV1}\par Special case: $\alpha=0$, $E=\mathbb{R}^d$
& $\mathcal{O}\!(1+\sqrt\kappa\,\ell)$\par $\varepsilon\gtrsim\sqrt\kappa\,\delta^2/\mu$
& i) \textbf{(1)}: certified threshold $\delta^2/\mu$ instead of $\sqrt\kappa\,\delta^2/\mu$ at $\alpha=0$.\newline ii) \textbf{(1)}: scope: permits $\alpha>0$.
& \textbf{(1)}: larger leading constant at $\alpha=0$. \\
\midrule
Gradient descent \citep[Thm.~5.2]{vasinGasnikov2026CompositeV1}
& $\mathcal{O}\!\left(1+\frac{\kappa}{(1-\alpha)^3}\ell\right)$\par $\varepsilon\gtrsim\frac{\delta^2}{\mu(1-\alpha)^3}$
& \textbf{(1)}: $\sqrt\kappa$ instead of $\kappa$ when $\alpha\le\kappa^{-1/4}$, $\kappa\gg1$, at accuracies certified by each bound.
& \textbf{(1)}: requires strong convexity instead of the PL condition. Comparison restricted to strongly convex $f$ on $\mathbb{R}^d$. \\
\bottomrule
\end{tabularx}
\begin{tablenotes}[flushleft]\footnotesize
\item[] \textbf{Abbr:} $\kappa=L/\mu$; $\ell=\log(1+LR^2/\varepsilon)$; one call returns one inexact gradient vector. PL: $\|\nabla f(x)\|^2\ge2\mu(f(x)-f(x_\star))$. At $\alpha=0$, $\alpha^{-1}=+\infty$. The symbol $\gtrsim$ in accuracy conditions requires a sufficiently large universal constant; $\asymp$ denotes equality up to universal factors.
\end{tablenotes}
\end{threeparttable}
\end{table}
\clearpage
